\documentclass{article}
\usepackage{mathtools} 
\usepackage{fontspec}
\usepackage[UTF8]{ctex}
\usepackage{amsthm}
\usepackage{mdframed}
\usepackage{xcolor}
\usepackage{amssymb}
\usepackage{amsmath}


% 定义新的带灰色背景的说明环境 zremark
\newmdtheoremenv[
  backgroundcolor=gray!10,
  % 边框与背景一致，边框线会消失
  linecolor=gray!10
]{zremark}{说明}

% 通用矩阵命令: \flexmatrix{矩阵名}{元素符号}{行数}{列数}
\newcommand{\flexmatrix}[4]{
  \[
  #1 = \begin{pmatrix}
    #2_{11}     & #2_{12}     & \cdots & #2_{1#4}   \\
    #2_{21}     & #2_{22}     & \cdots & #2_{2#4}   \\
    \vdots      & \vdots      & \ddots & \vdots     \\
    #2_{#31}    & #2_{#32}    & \cdots & #2_{#3#4}
  \end{pmatrix}
  \]
}

% 简化版命令（默认矩阵名为A，元素符号为a）: \quickmatrix{行数}{列数}
\newcommand{\quickmatrix}[2]{\flexmatrix{A}{a}{#1}{#2}}

\begin{document}
\title{习题 6.1}
\author{张志聪}
\maketitle

\section*{2}

 (2)

令$f(x) = x^n + px + q$
\begin{itemize}
  \item $n$是偶数。

        \begin{align*}
          f'(x) = n x^{n-1} + p \\
        \end{align*}
        于是
        \begin{align*}
          f'(x) = 0         \\
          n x^{n-1} + p = 0 \\
          x = \sqrt[n - 1]{-\frac{p}{n}}
        \end{align*}
        只有一个实数根。
        反证法，假设$f(x)$又两个不同的实数根，通过罗尔定理得知，
        存在两个$x_1, x_2$使得
        \begin{align*}
          f'(x_1) = f'(x_2) = 0
        \end{align*}
        出现矛盾。

  \item $n$是奇数。

        \begin{align*}
          f'(x) = n x^{n-1} + p \\
        \end{align*}
        于是，如果$p > 0$，那么$f'(x) = 0$没有根。

        如果$p = 0$，那么$f'(x) = 0$，只有一个实数根$x = 0$。

        如果$p < 0$，那么$f'(x) = 0$，只有两个实数根$x = \pm \sqrt[n - 1]{-\frac{p}{n}}$。

        所以，类似的推论，$f(x)$至多有三个不同的实数根。


\end{itemize}


\section*{10}

提示：推论3（导数极限定理）和定理3.10。


$f'(x)$在区间上单调，则$f'(x)$只有第一类间断点（第四章$\S 1$习题6）。

$f$在$(a, b)$上可导，则$f'(x)$没有第一类间断点（详细证明查看6-1-comnent.tex）。

综上，$f'(x)$在$(a, b)$上连续。

\section*{12}

提示：多次使用罗尔定理。

\section*{14}

因为$x \in (0, \frac{\pi}{2}), x sinx > 0$，
使用交叉相乘，令
\begin{align*}
  f(x) = tanx sinx - x^2
\end{align*}

我们只需证明$f(x) > 0$即可。




\end{document}